\documentclass[10pt,a4paper]{article}
\usepackage{amsmath,amssymb,amsthm,graphicx}
\usepackage{amsmath,amssymb,amsthm,hyperref,xfrac}
\providecommand{\ie}{i.e.\ }
\providecommand{\skpt}{\smallskip}
\usepackage[scr=rsfs,cal=euler]{mathalfa}
\multlinegap0pt

\hyphenation{begin decom-po-si-tion sequence}
\newenvironment{enumeratea}
{\bgroup\def\theenumi{\alph{enumi}}\begin{enumerate}}
{\end{enumerate}\egroup}

\newcommand{\ba}{{\boldsymbol a}}
\newcommand{\bb}{{\boldsymbol b}}
\newcommand{\bd}{{\boldsymbol d}}
\newcommand{\bu}{{\boldsymbol u}}
\newcommand{\bx}{{\boldsymbol x}}
\newcommand{\bmy}{{\boldsymbol y}}

\newcommand{\bD}{{\boldsymbol D}}
\newcommand{\bL}{{\boldsymbol L}}
\newcommand{\bP}{{\boldsymbol P}}
\newcommand{\bV}{{\boldsymbol V}}
\newcommand{\bX}{{\boldsymbol X}}
\newcommand{\bZ}{{\boldsymbol Z}}

\newcommand{\Psfrac}[2]{(\sfrac{#1}{#2})}
\newcommand{\psfrac}[2]{\sfrac{(#1)}{#2}}
\newcommand{\bigO}[1]{O(#1)}

\usepackage{mathtools}
\usepackage{bbm}

\newcommand{\NN}{\mathbb{N}}
\newcommand{\ZZ}{\mathbb{Z}}
\newcommand{\RR}{\mathbb{R}}
\newcommand{\CC}{\mathbb{C}}
\newcommand{\PP}{\mathbb{P}}
\newcommand{\EE}{\mathbb{E}}

\newcommand{\boldalpha}{\boldsymbol{\alpha}}
\newcommand{\bolddelta}{\boldsymbol{\delta}}
\newcommand{\boldrho}{\boldsymbol{\rho}}
\newcommand{\dee}{\,\mathrm{d}}
\newcommand{\one}{\mathbbm{1}}
\newcommand{\BC}{\mathcal{B}}
\newcommand{\EC}{\mathcal{E}}
\newcommand{\RS}{\mathscr{R}}
\newcommand{\prob}[1]{\PP[ #1 ]}
\newcommand{\bigoo}[1]{O\!\left(#1 \right)}
\newcommand{\abs}[1]{\left\lvert #1 \right\rvert}
\newcommand{\floor}[1]{\left\lfloor #1 \right\rfloor}
\newcommand{\riemint}[4]{\int_{#1}^{#2}\! #3 \, \mathrm{d}#4}
\newcommand{\maxnorm}[1]{\left\| #1 \right\|_{\infty}}

\renewcommand{\Re}{\textup{Re}}
\renewcommand{\Im}{\textup{Im}}

\newcommand{\dtv}{\mathrm{d}_{\mathrm{TV}}}
\newcommand{\dw}{\mathrm{d}_{\mathrm W}}

\let\epsilon\varepsilon

\newcommand{\pextra}{P_{\mathtt{extra}}}
\newcommand{\parratia}{P_{\mathtt{Arratia}}}

\newcommand{\simplex}{\Delta^{k-1}}

\newcommand{\bg}{\big}
\newcommand{\bgg}{\Big}
\newcommand{\bggg}{\bigg}

\newcommand{\Rnt}{R_{\mathrm{NT}}}
\newcommand{\Rprob}{R_{\mathrm{PD}}}
\newcommand{\ECnt}{\EC_{\mathrm{NT}}}
\newcommand{\ECprob}{\EC_{\mathrm{PD}}}

\newcommand{\Bnt}{\mathcal{B}_{\mathrm{NT}}}
\newcommand{\Bprob}{\mathcal{B}_{\mathrm{PD}}}

\DeclareMathOperator{\vard}{Var}
\DeclareMathOperator{\Dir}{Dir}

\theoremstyle{plain}
\newtheorem{cor}{Corollary}[section]
\newtheorem{lem}[cor]{Lemma}
\newtheorem{prop}[cor]{Proposition}
\newtheorem{conj}[cor]{Conjecture}
\newtheorem{theo}{Theorem}

\theoremstyle{definition}
\newtheorem{example}{Example}
\newtheorem{deff}[cor]{Definition}
\newtheorem{rem}{Remark}[section]
\newtheorem*{rems}{Remark}
\newtheorem*{remss}{Remarks}

\makeatletter
\ams@newcommand{\multiint}[1]{\DOTSI\protect\MultiIntegral{#1}}
\renewcommand{\MultiIntegral}[1]{%
\edef\ints@c{\noexpand\intop
\ifnum#1=\z@\noexpand\intdots@\else\noexpand\intkern@\fi
\replicate{#1-2}{\noexpand\intop\noexpand\intkern@}%
\noexpand\intop
\noexpand\ilimits@
}%
\futurelet\@let@token\ints@a
}
\makeatother
\ExplSyntaxOn
\cs_new:Npn \replicate #1 #2 { \prg_replicate:nn { #1 } { #2 } }
\ExplSyntaxOff
\title{Uniform fractional moments of the integer-friendly truncation N\_x: E[(x/N\_x)\textasciicircum{}alpha s(N\_x)\textasciicircum{}beta] is bounded}
\author{Tony Haddad, Dimitris Koukoulopoulos}
\date{doi:10.5802/jep.317; CC BY 4.0}
\begin{document}
\maketitle
\noindent\textit{Editorial note.} The counted statement is the author's labelled lemma `lem:properties of N` (source lines 517-523); the proof environment begins on the next line after \textbackslash{}end\{lem\} apart from one blank line, so the pairing is by adjacency.\par\medskip
\noindent\textit{Editorial note.} The retained context carries the author's own lead-in sentence, which recalls the definition of s(n) that the statement and the proof both use, together with the display defining Theta\_x.\par\medskip
\noindent\textit{Editorial note.} No counter adjustment is retained: the publisher numbering of the source is not reproduced, and the item prints with the wrapper's own sequential numbering.\par\medskip
\noindent\textit{Editorial note.} No retained line refers to a label, citation or figure outside the retained span, so the package declares no external reference.\par\medskip
\noindent\textit{Editorial scope.} One original result of the pinned publisher article, counted once: the uniform fractional-moment estimate for the integer-friendly truncation N\_x (author label lem:properties of N).\par\medskip
\noindent\textit{Allowed standard prerequisites (editorial).}\begin{itemize}
\item Square-full divisors. s(n) denotes the largest square-full divisor of n, and n = n\textasciicircum{}flat * s(n) with n\textasciicircum{}flat the square-free part.. The proof splits n as a * b with b = s(n) at source line 529 and then sums over square-full b.
\item Divisor-summation and partial sums of a power. Partial sums of a\textasciicircum{}\{-alpha\} for alpha < 1 grow like (x/b)\textasciicircum{}\{1-alpha\}.. It is the bound applied at source lines 531-532 to the inner sum over a.
\item Convergence of the square-full series. The series over square-full b of b\textasciicircum{}\{beta-1\} converges when beta < 1/2, since the square-full numbers have density of order x\textasciicircum{}\{1/2\}.. It is exactly the assumption beta < 1/2 used at source line 536 to close the estimate.
\end{itemize}\medskip
\begin{equation}
\label{eq:def of Theta_x}
\Theta_x \coloneqq \sum_{i \ge 1} r(V_i\log x).
\end{equation}

\begin{lem}[Properties of $N_x$]
\label{lem:properties of N}
Fix $\alpha\in[0,1)$ and $\beta\in[0,1/2)$. Uniformly over $x\ge1$, we~have
\[
\EE\bg[(x/N_x)^\alpha s(N_x)^\beta\bg] \ll_{\alpha,\beta} 1.
\]
\end{lem}

\begin{proof} We must show that
\[
S\coloneqq \sum_{n \le x} (x/n)^\alpha s(n)^\beta\ll_{\alpha,\beta} x.
\]
Indeed, if we let $b=s(n)$ and $a=n^\flat=n/b$, then
\[
S\le \sum_{\substack{b\le x \\ b\ \text{square-full} }} x^\alpha b^{\beta-\alpha} \sum_{a\le x/b} a^{-\alpha}
\ll_\alpha \sum_{\substack{b\le x \\ b\ \text{square-full} }} x^\alpha b^{\beta-\alpha} (x/b)^{1-\alpha}
\le x \sum_{b\ \text{square-full} } b^{\beta-1},
\]
where we used our assumption that $\alpha<1$.
Since we have assumed that $\beta<1/2$, the last sum over $b$ converges, thus completing the proof.
\end{proof}

\begin{thebibliography}{99}
\bibitem{Arratia98} R. Arratia - “On the central role of scale invariant Poisson processes on (0,∞)”, in Microsurveys in discrete probability (Princeton, NJ, 1997), DIMACS Ser. Discrete Math. Theoret. Comput. Sci., vol. 41, American Mathematical Society, Providence, RI, 1998, p. 21-41
\bibitem{Arratia02} R. Arratia - “On the amount of dependence in the prime factorization of a uniform random integer”, in Contemporary combinatorics, Bolyai Soc. Math. Stud., vol. 10, János Bolyai Math. Soc., Budapest, 2002, p. 29-91
\bibitem{ABT06} R. Arratia, A. D. Barbour \& S. Tavaré - “A tale of three couplings: Poisson-Dirichlet and GEM approximations for random permutations”, Combin. Probab. Comput. 15 (2006) no. 1-2, p. 31-62
\bibitem{BareikisManstavicius07} G. Bareikis \& E. Manstavičius - “On the DDT theorem”, Acta Arith. 126 (2007) no. 2, p. 155-168
\bibitem{Bill72} P. Billingsley - “On the distribution of large prime divisors”, Period. Math. Hungar. 2 (1972), p. 283-289
\bibitem{delaBretecheTenenbaum16} R. de la Bretèche \& G. Tenenbaum - “Sur les processus arithmétiques liés aux diviseurs”, Adv. in Appl. Probab. 48 (2016) no. A, p. 63-76
\bibitem{DDT79} J.-M. Deshouillers, F. Dress \& G. Tenenbaum - “Lois de répartition des diviseurs. I”, Acta Arith. 34 (1979) no. 4, p. 273-285 (loose errata)
\bibitem{DonnellyTavare87} P. Donnelly \& S. Tavaré - “The population genealogy of the infinitely-many neutral alleles model”, J. Math. Biol. 25 (1987) no. 4, p. 381-391
\bibitem{Feng10} S. Feng - The Poisson-Dirichlet distribution and related topics. Models and asymptotic behaviors, Probability and its Appl., Springer, Heidelberg, 2010
\bibitem{JKB} N. L. Johnson, S. Kotz \& N. Balakrishnan - Discrete multivariate distributions, Wiley Series in Probab. and Stat.: Applied Probab. and Stat., John Wiley \& Sons, Inc., New York, 1997
\bibitem{Kingman93} J. F. C. Kingman - Poisson processes, Oxford Studies in Probab., vol. 3, The Clarendon Press, Oxford University Press, New York, 1993
\bibitem{KotzBal00} S. Kotz, N. Balakrishnan \& N. L. Johnson - Continuous multivariate distributions. Vol. 1. Models and applications, Wiley Series in Probab. and Stat.: Applied Probab. and Stat., Wiley-Interscience, New York, 2000
\bibitem{Koukoulo19} D. Koukoulopoulos - The distribution of prime numbers, Graduate Studies in Math., vol. 203, American Mathematical Society, Providence, RI, 2019
\bibitem{LastPenrose18} G. Last \& M. Penrose - Lectures on the Poisson process, Institute of Math. Stat. Textbooks, vol. 7, Cambridge University Press, Cambridge, 2018
\bibitem{Ken22} S.-K. Leung - “Dirichlet law for factorisation of integers, polynomials and permutations”, Math. Proc. Cambridge Philos. Soc. 175 (2023) no. 3, p. 649-676
\bibitem{KenWeb24} S.-K. Leung - “Dirichlet law with arbitrary parameters for factorizations of integers” (2025), Retrieved on 24 August 2025 from https://sites.google.com/view/sunkaileung/miscellaneous
\bibitem{MV07} H. L. Montgomery \& R. C. Vaughan - Multiplicative number theory. I. Classical theory, Cambridge Studies in Advanced Math., vol. 97, Cambridge University Press, Cambridge, 2007
\bibitem{NyandwiSmati13} S. Nyandwi \& A. Smati - “Distribution laws of pairs of divisors”, Integers 13 (2013), article ID A13, 13 pages
\bibitem{Tenenbaum00} G. Tenenbaum - “A rate estimate in Billingsley’s theorem for the size distribution of large prime factors”, Q. J. Math. 51 (2000) no. 3, p. 385-403
\end{thebibliography}
\end{document}
